\documentclass[11pt,a4paper]{article}
\usepackage[T2A]{fontenc}
\usepackage{babel}
\babelprovide[import,main]{russian}
\usepackage[hmargin=27mm,vmargin=24mm]{geometry}

\begin{document}

\begin{center}
  {\Large\bfseries Футбольная команда}
\end{center}

\section*{Алгоритм}
\begin{itemize}
  \setlength{\itemsep}{0pt}
  \item Храним в векторе положительные эффективности игроков и их исходные номера.
    Сортируем по эффективности: $a_1 \le \cdots \le a_n$.
  \item Для каждой правой границы $r$ ищем самую левую границу $\ell$, при которой отрезок $[\ell,r]$ задаёт
    сплочённую команду.
    К текущей сумме добавляем $a_r$ и вычитаем эффективности выбывших игроков.
  \item Подходящая самая левая граница ищется комбинацией экспоненциального и двоичного поиска
    для маржинального ускорения.
  \item Копируем лучший отрезок в вектор ответа, сортируем его по исходным номерам и выводим сумму и номера.
  \item Обе сортировки используют быструю сортировку с трёхчастным разбиением и сортировку вставками на отрезках
    до 64 игроков.
\end{itemize}

\section*{Корректность}
\textbf{Лемма об отрезке.}
Команда из трёх и более игроков сплочённа, если и только если максимум не превосходит суммы двух минимумов.
Следует из того, что любая другая сумма гарантированно не меньше суммы двух минимумов.

\medskip
\noindent\textbf{Лемма о линейности поиска.}
В отсортированном массиве правая граница проходит $n$ позиций, а левая не убывает:
при расширении отрезка вправо максимум не уменьшается, а два минимума для отвергнутой границы не меняются.
Шаг со сдвигом левой границы на $d$ позиций стоит $O(1+d)$ (см. ниже).
Суммарный сдвиг не превосходит $n$, поэтому весь поиск занимает $O(n)$.
Поскольку эффективности положительны, самый левый допустимый старт даёт максимальную сумму для данного $r$.
Максимум по всем $r$ даёт оптимальную команду.

\section*{Время}
\begin{itemize}
  \setlength{\itemsep}{0pt}
  \item При сдвиге левой границы на $d$ позиций порог $\lceil a_r/2\rceil$ достигается не дальше чем через $d+1$ позиций:
    левая граница совпадает с первой позицией, достигшей порога, или стоит сразу перед ней.
    Поиск занимает $O(1+\log(d+1))$, обновление суммы --- $O(1+d)$.
    Весь проход занимает $O(n)$.
  \item Чтение, копирование, подсчёт суммы и вывод требуют $O(n)$ операций над числами.
  \item Сортировки вместе занимают ожидаемое $O(n\log n)$ при равномерном случайном выборе опорного элемента
    и $O(n^2)$ в худшем случае.\newline
    В коде используется \texttt{std::uniform\_int\_distribution} с \texttt{std::mt19937}.
\end{itemize}

\section*{Память}
\begin{itemize}
  \setlength{\itemsep}{0pt}
  \item Без векторов входа, ответа и копии входа --- $O(\log n)$: сортировка на месте рекурсивно
    обрабатывает меньшую часть, поиск занимает $O(1)$ памяти.
  \item Вход хранит $n$ игроков, ответ --- $k$. Вывод не требует дополнительного вектора.
  \item Вектор передаётся в решающую функцию по значению: результат чтения не копируется,
    а при передаче lvalue создаётся копия из $n$ элементов. Итого --- $O(n)$ в обоих случаях.
\end{itemize}

\end{document}
